\section{卷积神经网络概述}

\subsection{为什么需要卷积？}

到目前为止，我们唯一会用的网络架构是``全连接网络''——将图像展平为一维向量，然后做矩阵乘法。这有一个致命问题：\textbf{它破坏了图像的空间结构}。

\begin{figure}[H]
\centering
\includegraphics[width=0.95\textwidth]{slides-images/slide-019.jpg}
\caption{全连接层的缺陷：图像被展平后丢失了二维空间结构\protect\footnotemark}
\end{figure}
\footnotetext{Slides 第20页。}

图像是二维对象——像素之间的空间关系（邻近、对称、层次结构）对于理解图像内容至关重要。全连接层将图像视为无结构的数字序列，无法利用这些先验知识。

\begin{warningbox}{全连接网络处理图像的根本缺陷}
\begin{itemize}
    \item 展平操作破坏了二维空间结构
    \item 参数数量巨大（如 $224 \times 224 \times 3 = 150528$ 维输入，即使只有 1000 个隐藏单元也需要 1.5 亿参数）
    \item 无法利用图像的局部性和平移不变性
\end{itemize}
\end{warningbox}

\subsection{卷积网络的整体结构}

卷积神经网络（CNN / ConvNet）是由卷积层、池化层、全连接层和非线性激活函数构成的计算图，其一般结构为：

\begin{center}
\begin{tikzpicture}[node distance=0.6cm, every node/.style={draw, rounded corners, minimum width=2.2cm, minimum height=0.6cm, font=\small}]
\node (input) {输入图像};
\node (conv) [right=of input] {卷积 + ReLU};
\node (pool) [right=of conv] {池化};
\node (dots) [right=of pool, draw=none] {$\cdots$};
\node (fc) [right=of dots] {全连接层};
\node (out) [right=of fc] {分类分数};
\draw[->, thick] (input) -- (conv);
\draw[->, thick] (conv) -- (pool);
\draw[->, thick] (pool) -- (dots);
\draw[->, thick] (dots) -- (fc);
\draw[->, thick] (fc) -- (out);
\end{tikzpicture}
\end{center}

\begin{itemize}
    \item \textbf{前缀}（body）：交替的卷积层 + 激活函数 + 池化层，负责从原始像素中提取越来越抽象的特征表示
    \item \textbf{后缀}（head）：一个或多个全连接层，将特征映射到最终的预测分数
    \item 整个系统通过梯度下降端到端训练
\end{itemize}

\begin{figure}[H]
\centering
\includegraphics[width=0.95\textwidth]{slides-images/slide-021.jpg}
\caption{LeNet-5 架构（1998）：Yann LeCun 等人提出的经典卷积网络，用于手写数字识别\protect\footnotemark}
\end{figure}
\footnotetext{Slides 第22页。}

\begin{knowledgebox}{卷积网络的历史}
卷积神经网络的历史可以追溯到 1998 年 Yann LeCun 等人的 LeNet-5。尽管当时没有 GPU，计算资源极其有限，但网络架构的基本形式与 2010 年代的现代 CNN 几乎一致。2012 年 AlexNet 在 ImageNet 竞赛中的巨大成功开启了深度学习革命，从那时到约 2020 年，CNN 在计算机视觉的几乎每个任务上都占据统治地位。
\end{knowledgebox}

\subsection{CNN 的应用领域}

\begin{figure}[H]
\centering
\includegraphics[width=0.95\textwidth]{slides-images/slide-023.jpg}
\caption{CNN 的广泛应用：目标检测、语义分割、图像描述生成、文本到图像生成\protect\footnotemark}
\end{figure}
\footnotetext{Slides 第24页。}

从 2012 到 2020 年，CNN 几乎垄断了所有图像相关任务：
\begin{itemize}
    \item \textbf{目标检测}：在图像中画出每个物体的边界框并标注类别
    \item \textbf{语义分割}：为每个像素分配一个类别标签
    \item \textbf{图像描述生成}：从图像生成自然语言描述
    \item \textbf{图像生成}：Stable Diffusion 的第一个版本也基于卷积网络
\end{itemize}

\subsection{CNN vs Transformer}

\begin{figure}[H]
\centering
\includegraphics[width=0.85\textwidth]{slides-images/slide-025.jpg}
\caption{Vision Transformer（ViT, 2021）：将 Transformer 架构直接应用于图像\protect\footnotemark}
\end{figure}
\footnotetext{Slides 第26页。}

2017 年 Transformer 架构出现在 NLP 领域，2021 年 ViT（Vision Transformer）将几乎相同的架构应用于图像处理，并在许多任务上超越了 CNN。现在的趋势是 Transformer 逐渐取代 CNN 成为视觉任务的主流架构。

\begin{knowledgebox}{为什么还要学习 CNN？}
尽管 Transformer 在很多任务上超越了 CNN，学习 CNN 仍然非常重要：
\begin{enumerate}
    \item 巨大的\textbf{历史意义}——CNN 开启了深度学习时代
    \item CNN 在实践中\textbf{仍被广泛使用}
    \item 帮助建立\textbf{图像处理的直觉}
    \item CNN 并未完全``死亡''——很多现代系统是 CNN + Transformer 的\textbf{混合架构}
\end{enumerate}
\end{knowledgebox}

\subsection{本章小结}

卷积神经网络是专门为处理图像（以及其他具有空间结构的数据）设计的网络架构。它通过卷积层保留空间结构，通过池化层逐步缩小空间尺寸、扩大感受野，最终通过全连接层输出预测。尽管 Transformer 正在逐步取代 CNN，后者仍是理解现代视觉系统的重要基础。

%% ========== Part 4: 卷积层 ==========

